% Packages and macros for the thesis. The layout follows the thesis plan's template, with
% 1.5 line spacing (the EECS marksheet counts pages at 1.5-spaced 12 pt).

\usepackage[T1]{fontenc}
\usepackage[utf8]{inputenc}
\usepackage{lmodern}
\usepackage[margin=2.5cm]{geometry}
\usepackage{setspace}
\onehalfspacing
\usepackage{parskip}
\usepackage{microtype}

\usepackage{xcolor}
\usepackage{environ}   % the content notes' on/off switch (below)
\usepackage{graphicx}
\graphicspath{{figures/}}
\usepackage{array}
\usepackage{multirow}
\usepackage{makecell}
\usepackage{booktabs}
\usepackage{tabularx}
\usepackage{longtable}
\usepackage{caption}
\usepackage{subcaption}
\usepackage{float}

\usepackage{amsmath}
\usepackage{amssymb}

\usepackage{tikz}
\usetikzlibrary{positioning,arrows.meta,shapes.geometric,fit,calc,backgrounds}

\usepackage{url}
\urlstyle{tt}
\def\UrlBreaks{\do\-\do\/\do\.\do\_\do\?\do\&\do\=\do\:}
\usepackage[hidelinks]{hyperref}

\setlength{\emergencystretch}{3em}
\captionsetup{font=small,labelfont=bf}

\definecolor{instruc}{RGB}{112,48,160}
\definecolor{ubblue}{RGB}{10,132,255}
\definecolor{ubgrey}{RGB}{110,116,128}

% ---------------------------------------------------------------------------
% Content notes (the plan's own convention). Purple text is a note to the author:
% what the paragraph must say, the numbers it carries and where they come from.
%
% THE SWITCH: one line below.
%   \showinstructionstrue    every purple note is drawn (drafting)
%   \showinstructionsfalse   every note is hidden: the thesis as it reads without them, and for submission
% A clean copy beside the draft, without editing this file (run in docs/thesis):
%   latexmk -jobname=main-clean -usepretex="\def\hideinstructions{}" main.tex    -> build/main-clean.pdf
% A hidden note is dropped whole. A reference cited only inside notes leaves the bibliography until the
% prose cites it. Figure placeholders (\figplaceholder) stay visible either way.
% ---------------------------------------------------------------------------
\newif\ifshowinstructions
\showinstructionstrue
\ifdefined\hideinstructions\showinstructionsfalse\fi
\newcommand{\instructionstyle}{%
  \color{instruc}%
  \itshape%
  \fontsize{11pt}{13pt}\selectfont%
}
\newcommand{\instructiontext}[1]{\ifshowinstructions{\instructionstyle #1}\fi}
\NewEnviron{instructions}{%
  \ifshowinstructions\par\begingroup\instructionstyle\BODY\par\endgroup\fi}

% A source path inside a note: \src{docs/rq_roundA_2/RQ4b-detector-sweep.md}
\newcommand{\src}[1]{{\upshape\footnotesize\path{#1}}}

% Rotated table headings (AI-use table)
\newcommand{\rotatedheading}[1]{%
  \rotatebox[origin=c]{90}{%
    \makebox[4.7cm][l]{#1}%
  }%
}

% Placeholder for a figure not yet produced. Replace with \includegraphics once exported.
\newcommand{\figplaceholder}[2]{%
  \setlength{\fboxrule}{0.8pt}%
  \setlength{\fboxsep}{6pt}%
  \fcolorbox{instruc}{white}{%
    \parbox[c][#1][c]{0.88\linewidth}{\centering\instructionstyle #2}%
  }%
}

% The software's name, used the same way everywhere.
\newcommand{\tool}{MycoWave}
